\chapter{Routing methods and resource paths}
\chapterintro{The HTTP parser identifies a target; the router maps method and path shape to application code. Correct routing distinguishes absence from a method mismatch and makes precedence deterministic.}

\begin{objectives}
\begin{itemize}
  \item Register handlers by method and segment pattern.
  \item Match static and parameter segments with deterministic precedence.
  \item Distinguish 404, 405, automatic HEAD, and OPTIONS behavior.
  \item Keep query filtering outside path-shape matching.
\end{itemize}
\end{objectives}

\section{A small first router}

Stage 5 introduces a decorator, a route table, and \code{\{task\_id\}} capture:

\codefile{../code/stages/05_router/server.py}

The decorator records configuration at import time and returns the original
handler, so the function remains directly testable. Resolution splits paths into
segments and matches each against either literal text or a parameter placeholder.

The stage is intentionally naive: registration order can let a parameter route
shadow a static route, it reports every miss as 404, and it knows nothing about
HEAD or OPTIONS. The final router addresses each point.

\section{Static routes are more specific}

Consider \code{/users/me} and \code{/users/\{user\_id\}}. Both structurally match
the path \code{/users/me}, but the literal route expresses a more specific
resource. The final router sorts candidates by parameter count and then stable
registration order. Static routes therefore win without making behavior depend
on an accidental decorator order.

Only whole segments can be parameters. Patterns such as
\code{/files/\{name\}.json} are outside the small router's grammar. Parameter
names must be valid, unique Python identifiers. These rules keep matching easy
to teach and reason about.

\codefile{../code/final/tinyhttp/router.py}

\section{Path match and method match are separate facts}

Resolution gathers methods allowed by every route whose path pattern matches.
If no pattern matches, the result is 404. If a pattern matches but not for the
request method, the result is 405 with an \code{Allow} field. This distinction
helps clients discover and correct method use.

When GET is available, HEAD is also allowed. A HEAD request falls back to the GET
handler, after which the serializer suppresses content bytes. OPTIONS receives
an automatically generated 204 response with allowed methods. The application
does not require every resource handler to repeat protocol mechanics.

\section{Path parameters are strings first}

The router captures \code{42} from \code{/tasks/42} as a string. The Task layer
then decides that identifiers must be positive base-10 integers. This ownership
matters: a generic router should not assume every parameter is an integer, UUID,
slug, or date.

After a match, parameters are placed on the \code{Request}. A handler receives
one object rather than a changing positional signature. More advanced frameworks
offer typed converters and dependency injection; the explicit request structure
keeps this project's data flow visible.

\section{Queries refine a resource request}

The path chooses the collection \code{/tasks}. Query parameters such as
\code{status=todo}, \code{limit=20}, and \code{offset=0} refine the requested
representation. They do not participate in route selection. The parser provides
a mapping from each query name to a list of values so the application can reject
unexpected repetition instead of silently choosing one.

\begin{pitfall}
Do not route on a raw target string. A target can contain a query, percent
encoding, and repeated values. Parse the target once, match the decoded path by
well-defined segment rules, and validate query semantics at the resource layer.
\end{pitfall}

\section{Routing is not authorization}

Finding a handler means only that code exists for a method/path shape. A future
authorization layer would still decide whether the authenticated principal may
act on that particular resource. Returning 404 instead of 403 to hide existence
is an application security policy, not a router inevitability.

\summarytitle

The router maps method and decoded path segments to handlers, with static routes
ahead of parameter routes. It reports path absence separately from method
mismatch, supplies HEAD/OPTIONS behavior, and leaves parameter/query meaning to
the application.

\reviewtitle
\begin{enumerate}
  \item Why should \code{/users/me} win over \code{/users/\{user\_id\}}?
  \item What evidence distinguishes 404 from 405?
  \item Why does the router leave task ID conversion to the handler layer?
  \item Why should query parameters not be part of the route pattern here?
\end{enumerate}

\begin{exercise}[Route conflict table]
Register static and parameter routes for tasks, task statistics, and one task.
For ten paths and four methods, predict handler, captured parameters, allowed
methods, and final status. Then encode the table as tests.
\end{exercise}

\begin{exercise}[Add a nested resource]
Design patterns for \code{/projects/\{project\_id\}/tasks} and one nested task.
Identify which parameters the current router captures and which validation remains
in application code. Do not implement a database relation yet.
\end{exercise}
